\documentclass{article}

\usepackage{amsmath,amssymb}
\usepackage{xurl}
\usepackage[hidelinks]{hyperref}
\setlength{\emergencystretch}{2em}

\input{command}

\title{Top Schmidt-rank index in the $n$-positivity threshold of $T_\eta$
(Wolf Chapter 3)}
\author{}
\date{2026-06-24}

\begin{document}
\maketitle
\gapnote{scope-restriction}{resolved}

\section{Source statement}

Let $D\geq1$. Chapter~3 of Ref.~\cite{Wolf2012Quantum} studies the maps
$T_\eta:\MN{D}\to\MN{D}$ defined for $\eta\in\mathbb R_+$ by source
Equation~(3.11):
\begin{align}
  T_\eta(\rho)
    &=\tr(\rho)\Id-\frac{\rho}{\eta}.
  \label{eq:wolf_t_eta_top_index_scope_map}
\end{align}
The formal map is written $\tr(\rho)\Id-\eta^{-1}\rho$. It agrees with
(\ref{eq:wolf_t_eta_top_index_scope_map}) for $\eta>0$ and is also defined at
$\eta=0$.

The Choi--Jamiolkowski operator used in the source is
\begin{align}
  \tau_\eta
    &=D^{-1}\Id-\eta^{-1}\ketbra{\Omega}{\Omega}.
  \label{eq:wolf_t_eta_top_index_scope_choi_operator}
\end{align}
All positive eigenvalues of $\tau_\eta$ equal $1/D$. Its unique nonpositive
eigenvalue and the corresponding reduced density matrix are
\begin{align}
  \nu_-
    &=\frac{1}{D}-\frac{1}{\eta},
  \label{eq:wolf_t_eta_top_index_scope_negative_eigenvalue}\\
  \rho_-
    &=\frac{\Id}{D}.
  \label{eq:wolf_t_eta_top_index_scope_reduced_density}
\end{align}
Consequently, the Ky--Fan $n$-norm satisfies
\begin{align}
  \|\rho_-\|_{(n)}
    &=\frac{n}{D}.
  \label{eq:wolf_t_eta_top_index_scope_ky_fan_norm}
\end{align}
The source concludes that
\begin{align}
  T_\eta\text{ is $n$-positive}
    &\Longleftrightarrow \eta\geq n,
  \qquad 1\leq n\leq D.
  \label{eq:wolf_t_eta_top_index_scope_full_threshold}
\end{align}
At $n=D$, $n$-positivity is complete positivity. The threshold therefore makes
every inclusion in source Equation~(3.3) strict:
\begin{align}
  \mathcal T_{\mathrm{cp}}=\mathcal T_D
    &\subseteq\mathcal T_{D-1}
      \subseteq\cdots\subseteq\mathcal T_2
      \subseteq\mathcal T_1.
  \label{eq:wolf_t_eta_top_index_scope_positivity_chain}
\end{align}
These claims appear in Chapter~3 of Ref.~\cite{Wolf2012Quantum}.\footnote{Local
source:
\path{Notes/WolfNoteTexSource/ch03_positive_not_completely.tex:181--193},
equation label \path{eq:ch3:3.11}.}

\section{Point at issue}

The initial formal statement proved only
\begin{align}
  T_\eta\text{ is $k$-positive}
    &\Longleftrightarrow \eta\geq k,
  \qquad 1\leq k<D.
  \label{eq:wolf_t_eta_top_index_scope_initial_formal_threshold}
\end{align}
It omitted the top index $k=D$. This restriction was inherited from the formal
version of Wolf's maximal-overlap principle, Lemma~3.1
\cite[Lemma~3.1]{Wolf2012Quantum}. For normalized vectors $\psi$ of Schmidt
rank at most $k$, that principle identifies
\begin{align}
  \sup_\psi\left|\braket{\Omega}{\psi}\right|^2
    &=\|\rho\|_{(k)}.
  \label{eq:wolf_t_eta_top_index_scope_maximal_overlap}
\end{align}
The formal theorem was available only for $1\leq k<D$ because it depended on a
Ky--Fan maximum principle stated for orthogonal projections of rank exactly
$k<D$. At $k=D$, the Ky--Fan norm is the full trace and the maximizing
projection is $\Id$, but that boundary case was absent from the formal maximum
principle.

The resulting strictness witness, a $k$-positive map that is not
$(k+1)$-positive, was therefore available only when
\begin{align}
  1
    &\leq k,
  \label{eq:wolf_t_eta_top_index_scope_initial_witness_lower_bound}\\
  k+1
    &<D.
  \label{eq:wolf_t_eta_top_index_scope_initial_witness_upper_bound}
\end{align}
This range proves every strict inclusion in
(\ref{eq:wolf_t_eta_top_index_scope_positivity_chain}) except
\begin{align}
  \mathcal T_D
    &\subseteq\mathcal T_{D-1},
  \label{eq:wolf_t_eta_top_index_scope_top_inclusion}
\end{align}
which compares complete positivity with $(D-1)$-positivity and requires the
threshold at $k=D$.

\section{Resolution}

The endpoint $k=D$ is proved without the restricted maximal-overlap or Ky--Fan
principles. On $\MN{D}$, $D$-positivity is equivalent to complete positivity.
The corresponding formal equivalence is
\leanid{isCPMap_iff_isNPositiveMap_card}. Complete positivity of $T_\eta$ is
equivalent to positive semidefiniteness of
(\ref{eq:wolf_t_eta_top_index_scope_choi_operator}).

For every vector $\psi$, the quadratic form of the Choi operator is
\begin{align}
  \bra{\psi}\tau_\eta\ket{\psi}
    &=D^{-1}\|\psi\|^2
      -\eta^{-1}\left|\braket{\Omega}{\psi}\right|^2.
  \label{eq:wolf_t_eta_top_index_scope_choi_quadratic_form}
\end{align}
Because $\ket{\Omega}$ is normalized, Cauchy--Schwarz gives
\begin{align}
  \left|\braket{\Omega}{\psi}\right|^2
    &\leq\|\psi\|^2,
  \label{eq:wolf_t_eta_top_index_scope_cauchy_schwarz}
\end{align}
with equality at $\psi=\Omega$. Hence
\begin{align}
  \tau_\eta\succeq0
    &\Longleftrightarrow \eta\geq D,
  \label{eq:wolf_t_eta_top_index_scope_choi_threshold}\\
  T_\eta\text{ is $D$-positive}
    &\Longleftrightarrow \eta\geq D.
  \label{eq:wolf_t_eta_top_index_scope_top_threshold}
\end{align}

The declaration \leanid{Matrix.isNPositiveMap_tEta_iff} proves
(\ref{eq:wolf_t_eta_top_index_scope_initial_formal_threshold}), and
\leanid{Matrix.isNPositiveMap_tEta_card_iff} proves
(\ref{eq:wolf_t_eta_top_index_scope_top_threshold}). Both are in
\doclink{QICLean/Channel/NPositivityChainStrict}. Together they establish the
full source range in
(\ref{eq:wolf_t_eta_top_index_scope_full_threshold}).

The witness
\leanid{Matrix.exists_isNPositiveMap_not_succ} remains stated for $k+1<D$.
The endpoint thresholds at $k=D-1$ and $k=D$ nevertheless imply
\begin{align}
  \mathcal T_D
    &\subsetneq\mathcal T_{D-1},
  \label{eq:wolf_t_eta_top_index_scope_strict_top_inclusion}
\end{align}
although this strict inclusion is not recorded as a separate witness
declaration.

\section{Classification}

\emph{Resolved.} The declarations
\leanid{Matrix.isNPositiveMap_tEta_iff} and
\leanid{Matrix.isNPositiveMap_tEta_card_iff} now cover the complete source
range $1\leq k\leq D$ from source Equation~(3.11)
\cite[Chapter~3]{Wolf2012Quantum}. The top index uses the Choi-operator
argument rather than the restricted Ky--Fan route. No hypothesis foreign to
the source is introduced.

\bibliographystyle{alpha}
\bibliography{references}

\end{document}
